% !TeX root = ../main.tex
\section{Preliminaries}
\label{sec:preliminaries}

This section collects the notation and analytical ingredients used
throughout the paper. We first recall standard properties of the
primal function associated with smooth NC--SC minimax optimization.
We then formulate the inner maximization through an observable
gradient-residual interface and record quantitative bounds for the
primal-function estimators used below. Finally, we record a generic inexact trust-region estimate
that will be repeatedly specialized in the subsequent sections.


\subsection{Problem Setting and Standard Properties}
\label{sec:problem-setting}

We consider the smooth NC--SC minimax problem
\begin{equation}
    \min_{x\in\mathbb{R}^{d_x}} P(x),
    \qquad
    P(x):=\max_{y\in\mathbb{R}^{d_y}} f(x,y).
    \label{eq:ncsc-problem}
\end{equation}
For every $x$, define
\begin{equation}
    y^*(x):=\arg\max_{y\in\mathbb{R}^{d_y}} f(x,y).
    \label{eq:inner-maximizer}
\end{equation}
We write
\[
    P^*:=\inf_x P(x),
    \qquad
    \Delta_P:=P(x_0)-P^*.
\]
Throughout the paper, $\|\cdot\|$ denotes the Euclidean norm for
vectors and the spectral norm for matrices. For a symmetric matrix
$A$, $\lambda_{\min}(A)$ denotes its smallest eigenvalue.

\begin{assumption}[NC--SC regularity]
\label{ass:ncsc}
The payoff
$f:\mathbb{R}^{d_x}\times\mathbb{R}^{d_y}\to\mathbb{R}$
is twice continuously differentiable and satisfies:
\begin{enumerate}[label=(\roman*)]
    \item $\nabla f$ is $\ell$-Lipschitz;
    \item $\nabla^2 f$ is $\rho$-Lipschitz;
    \item $f(x,\cdot)$ is $\mu$-strongly concave for every $x$,
          where $0<\mu\le\ell$;
    \item $P^*>-\infty$.
\end{enumerate}
We write
\begin{equation}
    \kappa:=\frac{\ell}{\mu}.
    \label{eq:kappa}
\end{equation}
\end{assumption}

Under Assumption~\ref{ass:ncsc}, the maximizer $y^*(x)$ is unique and
$\nabla^2_{yy}f(x,y)$ is nonsingular. Define the Schur-complement map
\begin{equation}
\begin{aligned}
    \mathcal H(x,y)
    &:=
    \nabla^2_{xx}f(x,y)
    -
    \nabla^2_{xy}f(x,y)
    \bigl(\nabla^2_{yy}f(x,y)\bigr)^{-1}
    \nabla^2_{yx}f(x,y).
\end{aligned}
\label{eq:schur-definition}
\end{equation}

We use the following standard properties of the primal function.

\begin{lemma}[Primal-function geometry]
\label{lem:primal-geometry}
Under Assumption~\ref{ass:ncsc}, $y^*(\cdot)$ is
$\kappa$-Lipschitz and $P$ is twice continuously differentiable with
\begin{subequations}
\label{eq:primal-derivatives}
\begin{align}
    \nabla P(x)
    &=
    \nabla_x f(x,y^*(x)),
    \label{eq:primal-gradient}
    \\
    \nabla^2P(x)
    &=
    \mathcal H(x,y^*(x)).
    \label{eq:primal-hessian}
\end{align}
\end{subequations}
Moreover, $\nabla P$ is $(1+\kappa)\ell$-Lipschitz and
\begin{equation}
    \|\nabla^2P(x)-\nabla^2P(x')\|
    \le
    M_P\|x-x'\|,
    \qquad
    M_P:=4\sqrt{2}\,\kappa^3\rho.
    \label{eq:primal-hessian-lipschitz}
\end{equation}
\end{lemma}

These properties follow from
\citet[Lemmas~1--3]{luo2022finding}.

\begin{definition}[Approximate second-order stationarity]
\label{def:sosp}
For $\epsilon_g,\epsilon_H>0$, a point $x$ is an
$(\epsilon_g,\epsilon_H)$-second-order stationary point of $P$ if
\begin{equation}
    \|\nabla P(x)\|\le\epsilon_g,
    \qquad
    \lambda_{\min}(\nabla^2P(x))\ge-\epsilon_H.
    \label{eq:sosp}
\end{equation}
Our main target is an
$(\epsilon,\sqrt{M_P\epsilon})$-second-order stationary point.
\end{definition}


\subsection{Inner Residuals and Primal-Function Estimators}
\label{sec:corrected-estimators}

The distance $\|y-y^*(x)\|$ is useful analytically but is not directly
observable. We therefore formulate the inner maximization in terms of
the gradient residual $\|\nabla_yf(x,y)\|$.

\begin{lemma}[Residual--distance certificate]
\label{lem:inner-residual-certificate}
Under Assumption~\ref{ass:ncsc}, for every $x$ and $y$,
\begin{equation}
    \|y-y^*(x)\|
    \le
    \frac{\|\nabla_yf(x,y)\|}{\mu}
    \le
    \kappa\|y-y^*(x)\|.
    \label{eq:inner-residual-certificate}
\end{equation}
\end{lemma}

\begin{proof}
The first inequality follows from strong concavity and the
Cauchy--Schwarz inequality. The second follows from
$\ell$-smoothness and $\nabla_yf(x,y^*(x))=0$.
\end{proof}

We use the inner maximization only through the following interface.

\begin{definition}[Inner residual module]
\label{def:inner-residual-module}
Given an outer point $x$, a warm start $y^{\rm in}$, and a target
$\tau>0$, an inner residual module $\mathcal S$ returns
\begin{equation}
    y^{\rm out}
    =
    \operatorname{MaxRes}_{\mathcal S}
    (x,y^{\rm in},\tau)
    \quad\text{such that}\quad
    \|\nabla_yf(x,y^{\rm out})\|\le\tau.
    \label{eq:maxres-interface}
\end{equation}

We call $\mathcal S$ accelerated if there exist
$C_{\mathcal S}>0$ and $b_{\mathcal S}\ge1$, which may depend on the
implementation and on $\kappa$ but not on
$x$, $\tau$, the initial residual, $\epsilon$, or the call index,
such that a call with
\[
    R_{\rm in}:=
    \|\nabla_yf(x,y^{\rm in})\|
\]
uses at most
\begin{equation}
    C_{\mathcal S}\sqrt{\kappa}
    \left[
        b_{\mathcal S}
        +
        \left\lceil
        \log_2
        \max\left\{
            1,\frac{R_{\rm in}}{\tau}
        \right\}
        \right\rceil
    \right]
    \label{eq:inner-module-work}
\end{equation}
evaluations of $\nabla_yf$.
\end{definition}

The residual condition \eqref{eq:maxres-interface} is the only
inner-solver property needed for outer correctness. The work bound
\eqref{eq:inner-module-work} is used only when deriving total
inner-gradient complexity. Concrete realizations by accelerated
gradient descent and SCAR are introduced in
Sections~\ref{sec:known-parameter} and~\ref{sec:parameter-free},
respectively.

We count one evaluation of $\nabla_yf(x,y)$ as one inner-gradient
evaluation. Corrected-gradient and Schur-complement constructions,
trust-region subproblem solves, and linear algebra used to apply
$(\nabla^2_{yy}f)^{-1}$ are accounted for separately.

We use the following value, gradient, and Hessian approximations
at an approximate inner maximizer $y$. Define
\begin{equation}
    \widehat P(x,y)
    :=
    f(x,y)
    -
    \frac12
    \nabla_yf(x,y)^\top
    \bigl(\nabla^2_{yy}f(x,y)\bigr)^{-1}
    \nabla_yf(x,y),
    \label{eq:corrected-value}
\end{equation}
and
\begin{subequations}
\label{eq:corrected-derivatives}
\begin{align}
    \widehat g(x,y)
    &:=
    \nabla_xf(x,y)
    -
    \nabla^2_{xy}f(x,y)
    \bigl(\nabla^2_{yy}f(x,y)\bigr)^{-1}
    \nabla_yf(x,y),
    \label{eq:corrected-gradient}
    \\
    \widehat H(x,y)
    &:=
    \mathcal H(x,y).
    \label{eq:corrected-hessian}
\end{align}
\end{subequations}
Since $\nabla^2_{yy}f(x,y)\prec0$, the correction in
\eqref{eq:corrected-value} raises the raw payoff and is exact for a
concave quadratic inner problem.

The gradient in \eqref{eq:corrected-gradient} is the existing
implicit estimator studied by \citet{ablin2020superefficiency},
adapted to inner maximization. Its quadratic accuracy is established
in their Proposition~2.3. The same total-gradient expression and
the Schur-complement Hessian appear in \citet{zhang2020newton}.
We use these oracle ingredients and the value correction to obtain
the explicit bounds required by the complexity analysis.
MCN instead uses the raw plug-in gradient $\nabla_xf(x,y)$, while
both methods use the Schur-complement Hessian.

\begin{lemma}[Quantitative primal-function oracle bounds]
\label{lem:corrected-errors}
Under Assumption~\ref{ass:ncsc}, for every $x$ and $y$,
\begin{subequations}
\label{eq:corrected-error-bounds}
\begin{align}
    |\widehat P(x,y)-P(x)|
    &\le
    \frac{(3\kappa-1)\rho}{12}
    \|y-y^*(x)\|^3
    \le
    \frac{\kappa\rho}{4}
    \|y-y^*(x)\|^3,
    \label{eq:corrected-value-error}
    \\
    \|\widehat g(x,y)-\nabla P(x)\|
    &\le
    \frac{(1+\kappa)\rho}{2}
    \|y-y^*(x)\|^2,
    \label{eq:corrected-gradient-error}
    \\
    \|\widehat H(x,y)-\nabla^2P(x)\|
    &\le
    (1+\kappa)^2\rho
    \|y-y^*(x)\|.
    \label{eq:corrected-hessian-error}
\end{align}
\end{subequations}
Moreover, whenever
\begin{equation}
    \|y-y^*(x)\|\le\frac{\mu}{\rho},
    \label{eq:corrected-value-local-region}
\end{equation}
the corrected value satisfies
\begin{equation}
    |\widehat P(x,y)-P(x)|
    \le
    \frac{7\rho}{24}\|y-y^*(x)\|^3.
    \label{eq:corrected-value-local-error}
\end{equation}
\end{lemma}

The proof is given in
Appendix~\ref{app:corrected-estimators}.
For completeness, we recall the cancellation behind the known
quadratic gradient error in \eqref{eq:corrected-gradient-error}.
With
\[
    y_t:=y^*(x)+t\bigl(y-y^*(x)\bigr),
    \qquad t\in[0,1],
\]
the proof derives the exact identity
\begin{align}
    \widehat g(x,y)-\nabla P(x)
    &=
    \int_0^1
    \Big[
        \nabla^2_{xy}f(x,y_t)-\nabla^2_{xy}f(x,y)
        \notag\\
    &\qquad
        +
        \nabla^2_{xy}f(x,y)
        \bigl(\nabla^2_{yy}f(x,y)\bigr)^{-1}
        \bigl(
            \nabla^2_{yy}f(x,y)
            -
            \nabla^2_{yy}f(x,y_t)
        \bigr)
    \Big]
    \notag\\
    &\qquad\qquad\times
    \bigl(y-y^*(x)\bigr)\,dt.
    \label{eq:corrected-gradient-cancellation}
\end{align}
Hessian Lipschitz continuity then gives the quadratic bound in
\eqref{eq:corrected-gradient-error}.

For comparison, the raw plug-in gradient satisfies
\begin{equation}
    \|\nabla_xf(x,y)-\nabla P(x)\|
    \le
    \ell\|y-y^*(x)\|.
    \label{eq:raw-gradient-error}
\end{equation}
Under these upper bounds, an inner distance of order $O(\epsilon)$
is sufficient for an $O(\epsilon)$ gradient approximation with the
raw estimator, whereas an inner distance of order
$O(\sqrt{\epsilon})$ suffices for the implicit estimator, with the
regularity constants fixed.

Combining Lemmas~\ref{lem:inner-residual-certificate}
and~\ref{lem:corrected-errors} gives an observable form of the same
estimates.

\begin{corollary}[Residual-based estimator bounds]
\label{cor:residual-estimator-bounds}
Let
\[
    R_y(x,y):=\|\nabla_yf(x,y)\|.
\]
Under Assumption~\ref{ass:ncsc},
\begin{subequations}
\label{eq:residual-estimator-bounds}
\begin{align}
    |\widehat P(x,y)-P(x)|
    &\le
    \frac{(3\kappa-1)\rho}{12\mu^3}
    R_y(x,y)^3,
    \label{eq:residual-value-bound}
    \\
    \|\widehat g(x,y)-\nabla P(x)\|
    &\le
    \frac{(1+\kappa)\rho}{2\mu^2}
    R_y(x,y)^2,
    \label{eq:residual-gradient-bound}
    \\
    \|\widehat H(x,y)-\nabla^2P(x)\|
    &\le
    \frac{(1+\kappa)^2\rho}{\mu}
    R_y(x,y).
    \label{eq:residual-hessian-bound}
\end{align}
\end{subequations}
If
\begin{equation}
    R_y(x,y)\le\frac{\mu^2}{\rho},
    \label{eq:residual-local-region}
\end{equation}
then
\begin{equation}
    |\widehat P(x,y)-P(x)|
    \le
    \frac{7\rho}{24\mu^3}R_y(x,y)^3.
    \label{eq:residual-local-value-bound}
\end{equation}
\end{corollary}

\begin{proof}
By Lemma~\ref{lem:inner-residual-certificate},
$\|y-y^*(x)\|\le R_y(x,y)/\mu$.
Substitution into Lemma~\ref{lem:corrected-errors} proves the result.
\end{proof}


\subsection{Trust-Region Preliminaries}
\label{sec:tr-preliminaries}

We finally record the generic trust-region estimates used by the outer
methods. Let $F:\mathbb{R}^n\to\mathbb{R}$ be twice continuously
differentiable and suppose that its Hessian is $M$-Lipschitz:
\begin{equation}
    \|\nabla^2F(x)-\nabla^2F(z)\|
    \le M\|x-z\|.
    \label{eq:generic-hessian-lipschitz}
\end{equation}

The following Taylor estimates are standard.

\begin{lemma}[Second-order Taylor bounds]
\label{lem:taylor-bounds}
For all $x,d\in\mathbb{R}^n$,
\begin{subequations}
\label{eq:taylor-bounds}
\begin{align}
    \|\nabla F(x+d)-\nabla F(x)-\nabla^2F(x)d\|
    &\le
    \frac{M}{2}\|d\|^2,
    \label{eq:gradient-taylor}
    \\
    F(x+d)
    &\le
    F(x)+\nabla F(x)^\top d
    +\frac12d^\top\nabla^2F(x)d
    +\frac{M}{6}\|d\|^3.
    \label{eq:function-taylor}
\end{align}
\end{subequations}
\end{lemma}

These estimates are standard; see, e.g.,
\citet[Lemma~1.2.4]{nesterov2004introductory}.

We also use the classical characterization of a globally solved
quadratic trust-region subproblem.

\begin{lemma}[Global optimality for a quadratic trust-region problem]
\label{lem:tr-global-optimality}
Let $g\in\mathbb{R}^n$, let $B\in\mathbb{R}^{n\times n}$ be symmetric,
and let $R>0$. A vector $d$ is a global solution of
\begin{equation}
    \min_{\|u\|\le R}
    \quad
    g^\top u+\frac12u^\top Bu
    \label{eq:generic-tr-subproblem}
\end{equation}
if and only if there exists $\lambda\ge0$ such that
\begin{equation}
    (B+\lambda I)d=-g,
    \qquad
    B+\lambda I\succeq0,
    \qquad
    \|d\|\le R,
    \qquad
    \lambda(\|d\|-R)=0.
    \label{eq:tr-global-optimality}
\end{equation}
In particular, $\lambda>0$ implies $\|d\|=R$.
\end{lemma}

This characterization is classical; see, e.g.,
\citet[Section~8.7]{luenberger2021linear}.

The next estimate is the common one-step inequality underlying the
fixed-scale and adaptive UTR analyses as well as the parameter-free
trust-region constructions later in the paper.

\begin{lemma}[Common inexact trust-region step estimate]
\label{lem:common-inexact-tr-step}
Suppose \eqref{eq:generic-hessian-lipschitz} holds.
At a point $x$, let $g$ and a symmetric matrix $H$ satisfy
\begin{equation}
    \|g-\nabla F(x)\|\le\delta_0,
    \qquad
    \|H-\nabla^2F(x)\|\le\xi.
    \label{eq:common-current-errors}
\end{equation}
For $a\ge0$ and $R>0$, let $d$ be a global solution of
\begin{equation}
    \min_{\|u\|\le R}
    \quad
    q(u):=
    g^\top u+\frac12u^\top(H+aI)u,
    \label{eq:common-tr-model}
\end{equation}
let $\lambda$ be its trust-region multiplier, and write
$r:=\|d\|$.
Let a candidate gradient $g^+$ satisfy
\begin{equation}
    \|g^+-\nabla F(x+d)\|\le\delta_+.
    \label{eq:common-candidate-error}
\end{equation}
Then
\begin{subequations}
\label{eq:common-tr-estimates}
\begin{align}
    q(d)
    &\le
    -\frac{\lambda}{2}r^2,
    \label{eq:common-model-bound}
    \\
    F(x)-F(x+d)
    &\ge
    \frac{a-\xi}{2}r^2
    +\frac{\lambda}{2}r^2
    -\frac{M}{6}r^3
    -\delta_0r,
    \label{eq:common-function-bound}
    \\
    \|g^+\|
    &\le
    (a+\xi+\lambda)r
    +\frac{M}{2}r^2
    +\delta_0+\delta_+.
    \label{eq:common-gradient-bound}
\end{align}
\end{subequations}
\end{lemma}

The proof is deferred to
Appendix~\ref{app:common-tr-step}.
